\documentclass[11pt,a4paper]{article}
\usepackage[utf8]{inputenc}
\usepackage[T1]{fontenc}
\usepackage{lmodern}
\usepackage[french]{babel}
\usepackage{amsmath,amssymb,mathtools}
\usepackage{icomma}
\usepackage{xcolor}
\usepackage{tikz}
\usepackage{pgfplots}
\usepackage[most]{tcolorbox}
\usepackage{enumitem}
\usepackage{fancyhdr}
\usepackage[hidelinks]{hyperref}
\usepackage{titlesec}
\usepackage[a4paper,margin=2.2cm,headheight=26pt,headsep=10pt]{geometry}
\usetikzlibrary{arrows.meta,positioning,calc,intersections,angles,quotes}
\usepgfplotslibrary{fillbetween}
\pgfplotsset{compat=1.18}
\definecolor{primary}{RGB}{30,75,145}
\definecolor{primarydark}{RGB}{20,50,100}
\definecolor{defcol}{RGB}{0,114,178}
\definecolor{thmcol}{RGB}{0,135,110}
\definecolor{excol}{RGB}{0,130,85}
\definecolor{remarkcol}{RGB}{120,70,180}
\definecolor{attncol}{RGB}{200,40,55}
\definecolor{methcol}{RGB}{85,85,100}
\definecolor{areacol}{RGB}{120,160,230}
\definecolor{areacol2}{RGB}{230,150,80}
\newcommand{\dd}{\mathrm{d}}
\newcommand{\R}{\mathbb{R}}
\newcommand{\ua}{\,\text{u.a.}}
\newcommand{\tg}{\tan}\newcommand{\cotg}{\cot}
\tcbset{boxbase/.style={enhanced,breakable,boxrule=0.9pt,arc=2.5pt,left=3mm,right=3mm,top=2.3mm,bottom=2.3mm,fonttitle=\bfseries,coltitle=white,toptitle=0.6mm,bottomtitle=0.6mm}}
\newtcolorbox{methode}[1][]{boxbase,colback=methcol!7,colframe=methcol,sharp corners,title={\textsc{Comment faire}\ifx\relax#1\relax\relax\else\textmd{ : \itshape #1}\fi}}
\newtcolorbox{exemple}[1][]{boxbase,colback=excol!5,colframe=excol,title={\textsc{Exemple}\ifx\relax#1\relax\relax\else\textmd{ : \itshape #1}\fi}}
\newtcolorbox{idee}[1][]{boxbase,colback=defcol!5,colframe=defcol,title={\textsc{Idée clé}\ifx\relax#1\relax\relax\else\textmd{ : \itshape #1}\fi}}
\newtcolorbox{remarque}[1][]{boxbase,colback=remarkcol!6,colframe=remarkcol,title={\textsc{Remarque}\ifx\relax#1\relax\relax\else\textmd{ : \itshape #1}\fi}}
\newtcolorbox{attention}[1][]{boxbase,colback=attncol!6,colframe=attncol,title={\textsc{Attention}\ifx\relax#1\relax\relax\else\textmd{ : \itshape #1}\fi}}
\newtcolorbox{exo}[1][]{boxbase,colback={primary!4},colframe=primary,title={\textsc{Exercice #1}}}
\newtcolorbox{corr}[1][]{boxbase,colback={thmcol!5},colframe=thmcol,title={\textsc{Corrigé #1}}}
\tikzset{axe/.style={->,>=Stealth,black!65,line width=0.7pt},courbe/.style={primary,line width=1.3pt},courbe2/.style={areacol2!90!black,line width=1.3pt},dvert/.style={gray,dashed,line width=0.7pt}}
\pagestyle{fancy}\fancyhf{}
\fancyhead[L]{\small\textcolor{primarydark}{\textbf{Fiche 7} — Intégrale d'une fonction continue}}
\fancyhead[R]{\small\textcolor{primarydark}{\textsc{Thème 5 — Calculs d'aires}}}
\fancyfoot[C]{\small\thepage}
\renewcommand{\headrulewidth}{0.8pt}
\titleformat{\section}{\large\bfseries\color{primarydark}}{\thesection}{0.6em}{}

\begin{document}

\section*{Corrigé — Moyen : fonction négative et changement de signe}

\begin{corr}[M1]
Sur $[-2,2]$, $x^2-4\leqslant 0$ : la fonction est \emph{négative}, on prend la valeur absolue.
\[
\int_{-2}^{2}(x^2-4)\,\dd x=\left[\frac{x^3}{3}-4x\right]_{-2}^{2}
=\left(\frac{8}{3}-8\right)-\left(-\frac{8}{3}+8\right)=-\frac{16}{3}-\frac{16}{3}=-\frac{32}{3}.
\]
L'intégrale est négative ; l'aire est sa valeur absolue :
$\boxed{\ \mathcal{A}=\dfrac{32}{3}\ua\approx 10{,}67\ \text{u.a.}\ }$
\end{corr}

\begin{corr}[M2]
Sur $[0,3]$, $f(x)=-2x\leqslant 0$. Donc
\[
\int_0^3 (-2x)\,\dd x=\bigl[-x^2\bigr]_0^3=-9\quad\Longrightarrow\quad
\mathcal{A}=\left|-9\right|=\boxed{\,9\ua.\,}
\]
\end{corr}

\begin{corr}[M3]
\textbf{a)} $f(x)=x^2-2x=x(x-2)$, racines $0$ et $2$. Parabole tournée vers le haut : $f\leqslant 0$
sur $[0,2]$ et $f\geqslant 0$ sur $[2,3]$. Le signe change en $x=2$ : on \textbf{coupe}.

\textbf{b)} Primitive : $F(x)=\dfrac{x^3}{3}-x^2$.
\[
\int_0^2 f=\left[\frac{x^3}{3}-x^2\right]_0^2=\frac{8}{3}-4=-\frac{4}{3}\ \Rightarrow\ \mathcal{A}_1=\frac{4}{3};
\qquad
\int_2^3 f=\left[\frac{x^3}{3}-x^2\right]_2^3=0-\left(-\frac{4}{3}\right)=\frac{4}{3}=\mathcal{A}_2.
\]
\[
\mathcal{A}=\mathcal{A}_1+\mathcal{A}_2=\frac{4}{3}+\frac{4}{3}=\boxed{\,\frac{8}{3}\ua\approx 2{,}67\ \text{u.a.}}
\]

\smallskip
\begin{center}
\begin{tikzpicture}
  \begin{axis}[width=8.6cm,height=4.4cm,axis lines=middle,
      xmin=-0.5,xmax=3.4,ymin=-1.4,ymax=3.4,xtick={0,2,3},ytick=\empty,
      xlabel={$x$},ylabel={$y$},
      every axis x label/.style={at={(current axis.right of origin)},anchor=west},
      every axis y label/.style={at={(current axis.above origin)},anchor=south}]
    \addplot[domain=-0.3:3.2,samples=80,courbe]{x^2-2*x};
    \addplot[name path=F1,domain=0:2,samples=40,draw=none,forget plot]{x^2-2*x};
    \addplot[name path=F2,domain=2:3,samples=30,draw=none,forget plot]{x^2-2*x};
    \path[name path=z1] (axis cs:0,0)--(axis cs:2,0);
    \path[name path=z2] (axis cs:2,0)--(axis cs:3,0);
    \addplot[areacol2!45] fill between[of=F1 and z1];
    \addplot[areacol!45] fill between[of=F2 and z2];
    \node at (axis cs:1,-0.6){$\mathcal{A}_1$};
    \node at (axis cs:2.65,0.6){$\mathcal{A}_2$};
  \end{axis}
\end{tikzpicture}
\end{center}
\end{corr}

\begin{corr}[M4]
$f(x)=x^2-4$ s'annule en $x=2$ sur $[0,3]$ : $f\leqslant 0$ sur $[0,2]$, $f\geqslant 0$ sur $[2,3]$.
Primitive : $F(x)=\dfrac{x^3}{3}-4x$.
\[
\int_0^2 f=\frac{8}{3}-8=-\frac{16}{3}\ \Rightarrow\ \mathcal{A}_1=\frac{16}{3};
\qquad
\int_2^3 f=\bigl(9-12\bigr)-\left(\frac{8}{3}-8\right)=-3+\frac{16}{3}=\frac{7}{3}=\mathcal{A}_2.
\]
\[
\mathcal{A}=\frac{16}{3}+\frac{7}{3}=\boxed{\,\frac{23}{3}\ua\approx 7{,}67\ \text{u.a.}}
\]
\end{corr}

\begin{corr}[M5]
$f(x)=x^2-1$ a pour racines $\pm 1$. Sur $[-2,2]$ : $f\geqslant 0$ sur $[-2,-1]$, $f\leqslant 0$ sur
$[-1,1]$, $f\geqslant 0$ sur $[1,2]$. Primitive : $F(x)=\dfrac{x^3}{3}-x$.

\smallskip
\begin{center}
\begin{tikzpicture}
  \begin{axis}[width=9cm,height=4.4cm,axis lines=middle,
      xmin=-2.5,xmax=2.5,ymin=-1.6,ymax=3.4,xtick={-2,-1,1,2},ytick=\empty,
      xlabel={$x$},ylabel={$y$},
      every axis x label/.style={at={(current axis.right of origin)},anchor=west},
      every axis y label/.style={at={(current axis.above origin)},anchor=south}]
    \addplot[domain=-2.3:2.3,samples=90,courbe]{x^2-1};
    \addplot[name path=Fa,domain=-2:-1,samples=30,draw=none,forget plot]{x^2-1};
    \addplot[name path=Fb,domain=-1:1,samples=40,draw=none,forget plot]{x^2-1};
    \addplot[name path=Fc,domain=1:2,samples=30,draw=none,forget plot]{x^2-1};
    \path[name path=za] (axis cs:-2,0)--(axis cs:-1,0);
    \path[name path=zb] (axis cs:-1,0)--(axis cs:1,0);
    \path[name path=zc] (axis cs:1,0)--(axis cs:2,0);
    \addplot[areacol!45] fill between[of=Fa and za];
    \addplot[areacol2!45] fill between[of=Fb and zb];
    \addplot[areacol!45] fill between[of=Fc and zc];
  \end{axis}
\end{tikzpicture}
\end{center}

Par symétrie, les deux morceaux extérieurs ont la même aire :
\[
\int_1^2 f=\left[\frac{x^3}{3}-x\right]_1^2=\left(\frac{8}{3}-2\right)-\left(\frac{1}{3}-1\right)
=\frac{2}{3}+\frac{2}{3}=\frac{4}{3}\ \Rightarrow\ 2\times\frac{4}{3}=\frac{8}{3}.
\]
Morceau central :
\[
\int_{-1}^{1} f=\left[\frac{x^3}{3}-x\right]_{-1}^{1}=\left(\frac{1}{3}-1\right)-\left(-\frac{1}{3}+1\right)
=-\frac{2}{3}-\frac{2}{3}=-\frac{4}{3}\ \Rightarrow\ \frac{4}{3}.
\]
\[
\mathcal{A}=\frac{8}{3}+\frac{4}{3}=\frac{12}{3}=\boxed{\,4\ua.\,}
\]
\end{corr}

\end{document}
